\begin{lemma}[Scheide internal replacement-pair edges (10)]
Fix defective $\delta$.  If
$\alpha\ne\gamma(\delta)$ is missing at a vertex of
$Z_{\gamma(\delta)}$, then an individual edge identity of color $\alpha$
joins $u^1_\delta$ and $u^2_\delta$ under the original coloring.
\end{lemma}
\begin{proof}
The defective color $\delta$ is present throughout
$Z_{\gamma(\delta)}$, so the assumed $\alpha$ is neither $\delta$ nor the
excluded $\gamma(\delta)$.  Thus the switch from $\varphi$ to
$\varphi_\delta$ changes neither the missing witness in $Z$ nor any
$\alpha$-edge identity.  By
\noderef{ScheideReplacementMaximum} and
\noderef{ScheideGlobalMaximumOrder}, the replacement vertex set
$V(T_\delta)=Z\cup U_\delta$ is elementary and closed for
$\varphi_\delta$.

The original closed set $V(T)$ has no $\alpha$-edge from $Z$ to
$U_\delta$, since $U_\delta$ lies outside it.  In the closed replacement
tree, the unique $\alpha$-edge at either replacement vertex can neither
leave $V(T_\delta)$ nor end in $Z$.  Elementariness makes $\alpha$ present
at both replacement vertices, since it is already missing in $Z$.
Consequently their unique $\alpha$ incidences are one and the same edge
identity joining the two vertices.  Finally
\noderef{ScheideReplacementPairFixed} says every incident edge at this
pair has the same color under $\varphi$ and $\varphi_\delta$, so the edge
has color $\alpha$ in the original coloring as asserted.
\end{proof}
